% Colophon
% Plan: 0002

\chapter*{Colophon}
\addcontentsline{toc}{chapter}{Colophon}

\vspace{1cm}

\begin{description}
\item[Build hash:] \texttt{\GitShortHash}
\item[Full commit:] \texttt{\GitHash}
\item[Commit date:] \GitDate
\item[Build date:] \BuildDate
\item[TeX engine:] \directlua{tex.print(status.banner)}
\end{description}

\vspace{1cm}

This book was typeset using Lua\LaTeX{} and built with GNU Bash. The source is maintained under version control. Every build is reproducible from the commit hash above.

\subsection*{For researchers}

This PDF contains hidden annotations at every point where source material was imported. These annotations record the original filename, archive location, date, and a brief description. In most PDF viewers, you can access them through the annotations or comments panel. They are invisible during normal reading.

\subsection*{Design Principles}

This book is a single self-contained file. The entire work --- science,
narrative, illustrations, interactive navigation, glossary, citations ---
ships as one HTML document under one megabyte, or as one PDF. No external
dependencies. No images to fail to load. No server to go down. No paywall.
No app.

This is deliberate. The design targets:

\begin{description}
  \item[Signal\,:\,noise.] Every paragraph either advances a takeaway or
    blocks a failure mode. Content that does neither is cut.
  \item[Signal\,:\,bulk.] A five-thousand-word summary. The remaining two
    hundred fifty pages are the evidence. Read the summary. Dive into the
    evidence when you are curious for details.
  \item[Portability.] Readable on a phone, a laptop, or printed on paper.
    No special software. No login. The file is the book.
  \item[Survivability.] A single file can be copied, emailed, torrented,
    archived, and cached. It does not depend on any service continuing to
    exist.
  \item[Accessibility.] Semantic HTML. Designed for screen-reader
    compatibility. All illustrations are inline SVG with title attributes.
    No raster images.
\end{description}

These constraints are features, not limitations. A preparation document that
depends on a server, an app, or a publisher's continued goodwill is not a
preparation document. It is a hostage. If you are reading this on a device
that can display a web page, you have everything you need.

\vspace{1cm}

\textit{``It is easier to get forgiveness than permission.''}

\vspace{1cm}

\noindent Media and press inquiries: \texttt{energyscholar+press@gmail.com}

\noindent General correspondence: \texttt{energyscholar+hello@gmail.com}
